\section{The algebraic seed and its finite rational certificate}
\label{art:seed_certificate}

This appendix gives the objects and finite inequalities needed for
Lemma~\ref{art:seed_cycle}. All decimals defining a box are terminating
rationals. No rounded point is substituted for a root or an infinite orbit.

\subsection{Existence and uniqueness of the defining root}
Put
\begin{align*}
 U&=127c+16d,&V&=16c-127d-16,&R&=c-8d,\\
 K&=8c^2+8d^2+d-8,&T&=c^2-c+d^2,&Y&=(c-1)^2+d^2,
\end{align*}
and let
\begin{equation}\label{art:seed_polynomials}
 P(c,d)=
 \begin{pmatrix}264192RK-UV(U+129)\\16VTK-RUY\end{pmatrix}.
\end{equation}
The root is specified by $P(c,d)=0$ and
$\|(c,d)-\bar r\|_\infty\le10^{-30}$, where
\begin{align*}
 \bar r_1&=0.31848142762711863064645131250290703717912205023605,\\
 \bar r_2&=-1.02523211967574872548135087001300414284511909318444.
\end{align*}

\begin{lemma}[Rational root certificate]\label{art:root_certificate}
This closed box contains exactly one root. On the box all of
$U,V,R,K,T,Y,U+129$ are positive, and \eqref{art:seed_rho} holds.
\end{lemma}
\begin{proof}
Let $D=DP(\bar r)$ and define its inverse by the ordinary two-by-two
adjugate formula $A=\operatorname{adj}(D)/\det D$.
The following finite rational certificate uses only the coefficients of
\eqref{art:seed_polynomials} and the two displayed rationals:
\begin{equation}\label{art:root_arithmetic}
 \det D\ne0,\qquad
 \max_i |(AP(\bar r))_i|<10^{-50},\qquad
 \max_i\sum_{2\le|\nu|\le5}|a_{i,\nu}|\,|\nu|\le7000,
\end{equation}
where, for $\nu=(\nu_1,\nu_2)$,
\[
 a_{i,\nu}=-\frac{1}{\nu_1!\nu_2!}
     \sum_{j=1}^2 A_{ij}\,\partial^\nu P_j(\bar r).
\]
These formulas specify every rational in the certificate: the index set
contains the eighteen monomials of degrees two through five. The bounds
in \eqref{art:root_arithmetic} follow by substituting the rational center,
differentiating the displayed polynomials, and clearing positive
denominators in the three comparisons. Thus they are finite integer
inequalities, rather than tolerances for a numerically computed root.

For completeness, their analytic consequence follows directly.
The centered residual map $N(h)=h-AP(\bar r+h)$ has the exact expansion
\[
 N_i(h)=-(AP(\bar r))_i+
           \sum_{2\le|\nu|\le5}a_{i,\nu}h^\nu;
\]
its linear part vanishes because $AD=I$.
On $\|h\|_\infty,\|k\|_\infty\le r=10^{-30}$, telescoping the
factors of each monomial gives
$|h^\nu-k^\nu|\le|\nu|r\|h-k\|_\infty$, since $r\le1$ and
$|\nu|\ge2$. Hence $N$ has Lipschitz constant at most $7000r$ and
\[
 \|N(h)\|_\infty\le10^{-50}+7000r^2<r.
\]
The closed box is complete, so contraction gives a unique fixed point.
Since $A$ is invertible, this is exactly a root of $P$.
This is the standard preconditioned-residual principle; compare
Rump~\cite[Theorem~13.3]{NEG_RUMP10}.

The root box lies in
$318/1000<c<319/1000$ and $-1026/1000<d<-1025/1000$.
Substitution gives positivity of the stated factors. For example,
\[
 \frac{2397}{100}\le U\le\frac{24113}{1000},\quad
 \frac{119263}{1000}\le V\le\frac{59703}{500},\quad
 \frac{4259}{500}\le R\le\frac{8527}{1000}.
\]
Multiplying positive endpoints in $129R/[V(U+129)]$ gives
$60/1000<\rho<61/1000$. In particular all divisions used below are valid.
\end{proof}

\subsection{The matrix identity}
\begin{proof}[Proof of the algebraic part of Lemma~\ref{art:seed_cycle}]
The first update in \eqref{art:seed_matrices} is immediate from
$s_0=(-1,0)^T$ and $g_1-g_0=(-2,-1/8)^T$.
For arbitrary $(c,d)$ with the positive factors above, put
\[
 \beta=\frac{66048RK}{UV(U+129)},\quad
 \gamma=\frac{VTK}{RUY},\quad
 \widetilde\rho=\frac{2064TK}{U(U+129)Y}.
\]
The rational identities obtained from \eqref{art:failure_update} are
\begin{equation}\label{art:closure_identity}
 \mathcal B(\mathcal B(M_1,g_1,g_2,\beta),g_2,g_0,\gamma)
     =\widetilde\rho I,\qquad
 -\gamma\mathcal B(M_1,g_1,g_2,\beta)g_2
     =\widetilde\rho(g_0-g_2).
\end{equation}
Here is a way to check the identities without guessing matrix entries.
The middle update is $C+\beta E$, because its two projection factors
are independent of the step length and its rank-one term is linear in
that length. Requiring $(C+\beta E)g_2$ to be parallel to $g_0-g_2$
gives a linear equation for $\beta$; substitution of $M_1,g_1,g_2$
gives the displayed value. For a unit vector $w\perp(g_0-g_2)$,
substitution gives $w^T(C+\beta E)w=\widetilde\rho$.
With the displayed $\gamma$, the last step is
$\widetilde\rho(g_0-g_2)$. Its projection factors are the orthogonal
projector onto $(g_0-g_2)^\perp$, and its rank-one term is
$\widetilde\rho$ times the complementary projector. Their sum is
$\widetilde\rho I$, proving \eqref{art:closure_identity}.
Each substitution is a rational identity and can equivalently be
checked by clearing their nonzero denominators and expanding.

Equations~\eqref{art:seed_polynomials} give exactly
$\beta=1/4$ and $\gamma=1/16$. Moreover,
\[
 \widetilde\rho-\rho=
 \frac{2064R}{V(U+129)}\left(\gamma-\frac1{16}\right)=0.
\]
This gives the matrix and displacement closures. The formula for $z_0$
gives $s_0+s_1+s_2=-(1-\rho)z_0$.
The strict inequalities below imply
$(g_{j+1}-g_j)^Ts_j>0$ at each phase. Positive definiteness follows
inductively from the full update: its two nonnegative quadratic terms
cannot both vanish on a nonzero vector when $y^Ts>0$.
\end{proof}

\subsection{Sound enclosure of every query}
An interval $[a,b]$ encloses the corresponding real quantity. Addition,
negation, multiplication and reciprocal are evaluated by
\begin{align*}
 [a,b]+[c,d]&=[a+c,b+d],\qquad -[a,b]=[-b,-a],\\
 [a,b][c,d]&=[\min(ac,ad,bc,bd),\max(ac,ad,bc,bd)],
\end{align*}
and $[a,b]^{-1}=[1/b,1/a]$ when $0\notin[a,b]$.
These inclusion rules follow from order on the reals; induction on a
rational expression proves the soundness of every resulting bound.
All endpoints used here are rational, so no directed floating-point
rounding is needed.

For the query calculation it suffices to use the larger box
\begin{equation}\label{art:query_box}
 c\in\frac{[318481,318482]}{10^6},\quad
 d\in\frac{[-1025233,-1025232]}{10^6},\quad
 \rho\in\frac{[6019,6021]}{10^5}.
\end{equation}
The first two inclusions follow from the root box; substitution in
\eqref{art:seed_rho}, with positive denominator, gives the third.
Evaluate $z_j$ from \eqref{art:seed_points} and use
\[
 p_0=(-1,0)^T,\quad p_1=(127/256,1/16)^T,\quad
 p_2=16\rho(1-c,-d)^T.
\]
For a row with active index $i$, set $q=z_j+\alpha p_j$ and evaluate
the following rational expressions by the inclusion rules:
\[
 (g_i-g_k)^Tq (k\ne i),\quad -\alpha g_j^Tp_j,\quad
 \begin{cases}
 g_j^Tz_j+\alpha g_j^Tp_j/3-g_i^Tq,&\text{acceptance},\\
 g_i^Tq-g_j^Tz_j,&\text{rejection}.
 \end{cases}
\]
For an acceptance also evaluate $\alpha g_i^Tp_j$.
Every expression in its row is strictly greater than the following
positive rational lower bound, in the order of Table~\ref{art:trial_table}:
\begin{equation}\label{art:query_lower_bounds}
 \frac1{10^4}(729, 2417, 145, 75, 5964, 2682, 1041, 221, 20).
\end{equation}
The same interval operations give the value enclosures in that table.
For the three base points, the designated active gaps are positive as
well. Evaluating $3\pm g_i^Tq$ for every $i$, $2-q_2$, and the
designated active value $g_i^Tq$ gives positive lower bounds for all base
points and all nine queries. This proves
\eqref{art:seed_query_bounds}. All inputs, expressions, interval rules
and output bounds of this finite certificate are now specified. The
strict active gaps also give differentiability on a neighborhood of
each query. Together with the matrix identities, these calculations
complete the proof of Lemma~\ref{art:seed_cycle}.
